\par\noindent\textbf{(1)} 两区波速 \(c_j=\sqrt{T/\mu_j}\)、波数 \(k_j=\omega/c_j\)，横波力对速度的特性阻抗 \(Z_j=T/c_j=\sqrt{T\mu_j}\)。在界面设 \(y_1=(Ae^{ik_1x}+Be^{-ik_1x})e^{-i\omega t}\)、\(y_2=Ce^{ik_2x-i\omega t}\)。
\par\noindent\textbf{(2)} 节点无质量且张力相同，边界 \(A+B=C\) 及 \(k_1(A-B)=k_2C\)。故振幅反射系数 \(r=B/A=(k_1-k_2)/(k_1+k_2)=(Z_1-Z_2)/(Z_1+Z_2)\)，透射振幅 \(t=C/A=2Z_1/(Z_1+Z_2)\)。
\par\noindent\textbf{(3)} 弦波的平均功率通量 \(\langle P\rangle=(\omega^2/2)Z|A|^2\)。因此能量反射率 \(R=|r|^2\)，透射率 \(\mathcal T=(Z_2/Z_1)|t|^2=4Z_1Z_2/(Z_1+Z_2)^2\)。两者和为1，符合无质量无耗散节点的功率守恒。
\par\noindent\textbf{(4)} 若 \(Z_2\gg Z_1\)，位移反射系数 \(r\to-1\)，界面近似固定端，反射波相位翻转；若 \(Z_2\ll Z_1\)，\(r\to+1\)，近似自由端反射。有限界面质量或弯曲刚度会引入频率相关散射，不再仅由两侧阻抗决定。
